\chapter{Khám phá dịch vụ: Eureka}
\label{ch:eureka}

\section{Vấn đề}

Gateway phải chuyển tiếp \code{/api/courses} tới course-service. Tới địa
chỉ nào? Trong Compose nó có thể ghi cứng \code{course-service:8082} ---
nhưng khi đó mỗi replica mới, mỗi service dời chỗ, mỗi môi trường đều là
một thay đổi cấu hình. Khám phá dịch vụ (service discovery) thay địa chỉ
bằng \emph{tên được phân giải lúc chạy}: các service tự đăng ký vào một
registry; client hỏi registry.

\section{Eureka hoạt động thế nào}

Eureka (từ Netflix, tổ tiên của stack Spring Cloud) là một REST service
giữ một registry các instance:

\begin{itemize}
  \item \textbf{Đăng ký.} Lúc khởi động mỗi client POST danh tính của
  mình: tên ứng dụng (\code{spring.application.name}), host, cổng, URL
  health.
  \item \textbf{Heartbeat.} Mỗi 30\,s (mặc định) client gia hạn lease của
  mình. Bỏ lỡ vài lần gia hạn và instance bị loại bỏ (evict).
  \item \textbf{Cache phía client.} Bên tiêu thụ lấy registry về và cache
  lại (mặc định làm mới mỗi 30\,s), rồi cân bằng tải \emph{cục bộ} giữa
  các instance mà nó biết. Registry sập trong chốc lát không làm dừng
  lưu lượng --- bên gọi tiếp tục dùng cache của mình.
  \item \textbf{Tự bảo toàn.} Nếu quá nhiều heartbeat biến mất cùng lúc,
  Eureka giả định đó là sự cố mạng chứ không phải chết hàng loạt và ngừng
  loại bỏ. Ít bất ngờ hơn khi mạng chập chờn; dữ liệu cũ hơn khi sự cố
  là thật.
\end{itemize}

Phía tiêu thụ trong dự án này là URI \code{lb://course-service} của
gateway (Chương~4): ``tra tên này trong registry, chọn một instance.''

\begin{figure}[h]
\centering
\begin{tikzpicture}[node distance=5mm and 12mm]
  \node[boxdark] (eureka) {Eureka\\\scriptsize discovery-server:8761};
  \node[box, below left=of eureka] (a) {course-service\\\scriptsize phiên bản 1};
  \node[box, below=of eureka] (b) {course-service\\\scriptsize phiên bản 2};
  \node[box, below right=of eureka] (gw) {api-gateway};
  \draw[flow] (a) -- node[lbl,left]{đăng ký + heartbeat} (eureka);
  \draw[flow] (b) -- (eureka);
  \draw[flowdash] (gw) -- node[lbl,right]{lấy registry} (eureka);
  \draw[flow] (gw) -- node[lbl,below]{\scriptsize lb:// chọn một} (b);
\end{tikzpicture}
\caption{Đăng ký, heartbeat, và cân bằng tải phía client.}
\end{figure}

\section{Server, và thiết lập sống còn duy nhất trên k8s}

Module discovery-server là Eureka với gần như không có cấu hình --- một
annotation \code{@EnableEurekaServer} và một cổng. Cấu hình thú vị nằm
trên mỗi \emph{client}, và một dòng trong đó xứng đáng có mục riêng vì nó
quyết định cả hệ thống có chạy được trên Kubernetes hay không:

\begin{lstlisting}[language=yaml]
- name: EUREKA_INSTANCE_PREFER_IP_ADDRESS
  value: "true"
\end{lstlisting}

Mặc định một client đăng ký \emph{hostname} của nó. Trên laptop của bạn
thế là ổn. Bên trong một pod Kubernetes, hostname là \emph{tên pod} ---
\code{course-service-7d9f6c-x2vlq} --- thứ các pod khác không phân giải
được. Mọi đăng ký sẽ là một địa chỉ không ai gọi tới được. \emph{IP} của
pod, ngược lại, định tuyến được trên toàn cụm. Vì thế mọi Deployment
trong \code{deploy/k8s/base/} đều đặt cờ này: đăng ký IP, không phải tên.

\begin{pitfall}[mọi thứ đều đăng ký, không gì gọi được gì]
Triệu chứng: dashboard Eureka hiện mọi service đều UP, nhưng mọi request
đi qua gateway đều thất bại với \code{UnknownHostException}. Nguyên
nhân: hostname được đăng ký là tên pod (xem ở trên). Đây là lỗi
Eureka-trên-Kubernetes phổ biến nhất, và nó vô hình cho đến lần gọi
chéo service đầu tiên.
\end{pitfall}

\begin{decision}[vì sao vẫn dùng Eureka, khi Kubernetes đã có Service]
Kubernetes đã cho mỗi Service một tên DNS (\code{course-service:8082})
với cân bằng tải tích hợp sẵn --- khám phá dịch vụ là bản địa của nền
tảng. Eureka ở đây là sự dư thừa có chủ đích vì CV: stack cổ điển vẫn
phổ biến trong ngành, và việc di cư \emph{khỏi} nó (Chương~27: thay
\code{lb://} bằng DNS Service thuần, xóa hai đơn vị triển khai) tự thân
là loại công việc hiện đại hóa mà doanh nghiệp trả tiền để có. Xây cả hai
stack là để học cuộc di cư đó.
\end{decision}

\section{Bên ngoài dự án}

\emph{Consul} (HashiCorp) là lựa chọn thay thế chính ngoài JVM: dựa trên
agent, không phụ thuộc ngôn ngữ, với health check do agent thực hiện thay
vì heartbeat tự báo cáo, cộng một kho key-value có thể thay luôn cả
config-server. \emph{Khám phá dịch vụ bản địa của Kubernetes} là cực còn
lại: không có tiến trình registry nào cả --- các controller của nền tảng
theo dõi readiness của pod và cập nhật endpoint của Service; ``đăng ký''
đơn giản là một pod trở thành Ready (Chương~17 cho thấy bộ máy đó). Hãy
để ý mẫu hình: nền tảng cung cấp càng nhiều, ứng dụng càng phải mang ít.

\begin{hardware}
Các giá trị mặc định bảo thủ của Eureka (heartbeat 30\,s, cache client
30\,s, loại bỏ sau 90\,s) nghĩa là một instance đã crash vẫn có thể nhận
lưu lượng tới một phút. Với một replica cho mỗi service thì điều đó hầu
như không quan trọng --- bên gọi đằng nào cũng thất bại --- nhưng nó giải
thích vì sao các đợt triển khai cuốn chiếu ở Chương~23 chờ readiness của
Kubernetes, chứ không chờ trạng thái Eureka.
\end{hardware}

\begin{quiz}
\begin{enumerate}
  \item Lần theo những gì xảy ra từ lúc course-service khởi động đến lúc
  gateway định tuyến tới nó thành công: những lần đăng ký, lấy về, và
  cache nào có liên quan, và độ trễ trong trường hợp xấu nhất là bao
  nhiêu?
  \item Vì sao \code{prefer-ip-address} quan trọng trên Kubernetes nhưng
  không quan trọng trên Compose? (Mỗi nền tảng làm cho cái gì phân giải
  được?)
  \item Chế độ tự bảo toàn đánh đổi tính chất nào lấy tính chất nào? Cho
  một kịch bản nó giúp ích và một kịch bản nó gây hại.
  \item Nếu bỏ Eureka và trỏ gateway thẳng tới
  \code{http://course-service:8082}, bộ máy Kubernetes nào thay thế từng
  tính năng của Eureka?
\end{enumerate}
\end{quiz}
